\documentclass[11pt,a4paper]{article}
\usepackage[utf8]{inputenc}
\usepackage[T1]{fontenc}
\usepackage{amsmath, amssymb, amsthm}
\usepackage{geometry}
\usepackage{booktabs}
\usepackage{array}
\usepackage{parskip}
\usepackage{enumitem}

\geometry{margin=2.5cm}
\newcolumntype{L}[1]{>{\raggedright\arraybackslash}p{#1}}

\title{From the Discrete Objective to the Integral,\\and What Is Left To Do With It}
\author{Working notes: step 2 of the Schr\"odingerised Hamilton--Jacobi project}
\date{}

\begin{document}
\maketitle

\noindent\textbf{Purpose.} These notes record (i) exactly how the discrete
Almgren--Chriss objective becomes a continuous-time integral, and (ii) what
remains to be done with that integral to obtain the time-dependent
(Hamilton--Jacobi) equation. Nothing in Part~II has been checked numerically yet.

% =====================================================================
\section*{Part I --- How we arrive at the integral}
% =====================================================================

\subsection*{I.0 Starting point}

The multi-asset objective is Eq.~(15) of Almgren and Chriss applied to their
Eqs.~(29)--(30) (Sec.~2.1, p.~13; Appendix~A.1, pp.~36--37; page numbers read
from running headers, to be verified against the PDF):

\begin{equation}
  U[\mathbf{x}]
  = \sum_{j=1}^{M}\left[
      \tfrac{1}{2}\gamma_j X_j^2
      + \varepsilon_j \sum_{k=1}^{N}|n_{jk}|
      + \frac{\tilde{\eta}_j}{\tau}\sum_{k=1}^{N} n_{jk}^2
    \right]
  + \lambda \sum_{k=1}^{N} \tau\,\mathbf{x}_k^\top C\,\mathbf{x}_k
  \label{eq:disc}
\end{equation}

The asset-by-asset form assumes diagonal permanent- and temporary-impact
matrices (the diagonal case, Eq.~32 of the reference).

\subsection*{I.1 Separate the fixed offset from the schedule-dependent cost}

\begin{itemize}[leftmargin=1.5em]
  \item $\tfrac12\gamma_jX_j^2$ is fixed by the number of shares to liquidate.
  \item $\varepsilon_j\sum_k|n_{jk}|=\varepsilon_jX_j$ provided every trade is a
        sale ($n_{jk}\ge0$, no buying back).
\end{itemize}

For sales-only trading with exact liquidation, define the fixed offset
\[
  K(\mathbf X)=\sum_{j=1}^{M}\left(\tfrac12\gamma_jX_j^2+\varepsilon_jX_j\right).
\]
The full discrete objective is this offset plus the schedule-dependent terms.
The spread part is not constant if a finite terminal penalty allows residual
inventory. For the exact-liquidation problem, the reduced objective is:

\begin{equation}
  U'[\mathbf{x}] = \sum_{j=1}^{M}\frac{\tilde{\eta}_j}{\tau}\sum_{k=1}^{N} n_{jk}^2
  + \lambda \sum_{k=1}^{N} \tau\,\mathbf{x}_k^\top C\,\mathbf{x}_k
\end{equation}

\subsection*{I.2 Replace shares-per-interval by a selling rate}

Shares sold in an interval equal rate times interval length:
\begin{equation}
  n_{jk}=\tau\,v_j(t_k)
  \qquad\Longrightarrow\qquad
  \frac{\tilde\eta_j}{\tau}\,n_{jk}^2=\frac{\tilde\eta_j}{\tau}\,\tau^2 v_j^2
  =\tilde\eta_j\,\tau\,v_j(t_k)^2 .
\end{equation}
Both terms of $U'$ now have the form $\tau\times(\text{something at }t_k)$,
summed over $k$. The $1/\tau$ in the original coefficient is what produces this.

\subsection*{I.3 The sum over time becomes an integral}

A sum $\sum_k \tau f(t_k)$ is a Riemann sum, and
\begin{equation}
  \sum_{k=1}^{N}\tau\,f(t_k)\;\xrightarrow[\tau\to0]{}\;\int_0^T f(t)\,dt .
\end{equation}
This is the only reason the sum over $k$ disappears: it is the same quantity
written for infinitely many infinitesimal intervals.

\subsection*{I.4 The sum over assets becomes a matrix product}

The sum over $j$ is not removed; it is written compactly. With
$H=\mathrm{diag}(\tilde\eta_1,\dots,\tilde\eta_M)$:
\begin{equation}
  \sum_{j=1}^{M}\tilde\eta_j\,v_j^2=\mathbf{v}^\top H\,\mathbf{v},
  \qquad
  \mathbf{x}^\top C\,\mathbf{x}=\sum_{a,b=1}^{M}C_{ab}\,x_a x_b .
\end{equation}

\subsection*{I.5 The inventory update becomes a derivative}

$\mathbf{x}_k=\mathbf{x}_{k-1}-\mathbf{n}_k$, divided by $\tau$, gives
$\dot{\mathbf{x}}=-\mathbf{v}$. The boundary conditions are the full position at
the start and a flat position at the horizon.

\subsection*{I.6 Result}

For exact liquidation and sales-only trading, the continuum total cost is
\begin{equation}
  J_{\mathrm{full}}[\mathbf v]=K(\mathbf X)+J[\mathbf v].
\end{equation}
The fixed offset is retained in the total cost; only the schedule-dependent
integral $J$ is passed to the HJB and Riccati derivation below. The spread
offset is not fixed when the terminal inventory is allowed to vary.

\begin{equation}
  \boxed{\;
  J[\mathbf{v}]
  = \int_0^T \Big( \mathbf{v}(t)^\top H\,\mathbf{v}(t)
      + \lambda\,\mathbf{x}(t)^\top C\,\mathbf{x}(t) \Big)\,dt,
  \quad \dot{\mathbf{x}}=-\mathbf{v},\quad
  \mathbf{x}(0)=\mathbf{X},\quad \mathbf{x}(T)=\mathbf{0}\;}
  \label{eq:J}
\end{equation}

\begin{center}
\begin{tabular}{@{}L{3.8cm}L{10cm}@{}}
\toprule
\textbf{Term} & \textbf{Meaning} \\
\midrule
$\mathbf{v}(t)$ & Selling rate in each asset; the control \\
$\mathbf{x}(t)$ & Inventory vector; the state \\
$H$ & $\mathrm{diag}(\tilde\eta_j)$; in the limit $\tau\to0$, $\tilde\eta_j=\eta_j-\tfrac12\gamma_j\tau\to\eta_j$ \\
$\mathbf{v}^\top H\mathbf{v}$ & Slippage cost rate; penalises selling fast \\
$\lambda\,\mathbf{x}^\top C\mathbf{x}$ & Risk cost rate; penalises holding a large or concentrated position \\
\bottomrule
\end{tabular}
\end{center}

\textbf{Source.} Eq.~\eqref{eq:J} is derived here; it is not in Almgren and
Chriss, who work in discrete time and solve by calculus of variations.

% =====================================================================
\section*{Part II --- What is left to do with the integral}
% =====================================================================

The reduced integral~\eqref{eq:J} defines the schedule optimization for one
starting position. Under exact liquidation, add $K(\mathbf X)$ to its optimized
value to recover the full initial cost.
To get an equation that evolves in time, the cost must become a function of the
current state and the current time.

\subsection*{II.1 Define the cost-to-go}

\begin{equation}
  \mathcal{V}(\mathbf{x},t)=\min_{\mathbf{v}(\cdot)}\int_t^T
  \Big(\mathbf{v}^\top H\mathbf{v}+\lambda\,\mathbf{x}^\top C\mathbf{x}\Big)\,ds,
  \qquad \mathbf{x}(t)=\mathbf{x}.
\end{equation}
The cheapest possible remaining cost if you hold $\mathbf{x}$ at time $t$.

\subsection*{II.2 Apply Bellman's principle}

Split the horizon into the next short step $dt$ and the rest:
\begin{equation}
  \mathcal{V}(\mathbf{x},t)=\min_{\mathbf{v}}\Big[\big(\mathbf{v}^\top H\mathbf{v}
  +\lambda\,\mathbf{x}^\top C\mathbf{x}\big)dt+\mathcal{V}(\mathbf{x}-\mathbf{v}\,dt,\;t+dt)\Big].
\end{equation}
Expanding the last term to first order,
$\mathcal{V}(\mathbf{x}-\mathbf{v}\,dt,t+dt)\approx\mathcal{V}+\partial_t\mathcal{V}\,dt-\nabla\mathcal{V}\cdot\mathbf{v}\,dt$,
and cancelling $\mathcal{V}$ gives
\begin{equation}
  -\partial_t\mathcal{V}=\min_{\mathbf{v}}\Big[\mathbf{v}^\top H\mathbf{v}
  +\lambda\,\mathbf{x}^\top C\mathbf{x}-\mathbf{v}^\top\nabla_{\mathbf{x}}\mathcal{V}\Big].
\end{equation}

\subsection*{II.3 Minimise over the control}

The bracket is quadratic in $\mathbf{v}$. Setting its derivative to zero:
\begin{equation}
  2H\mathbf{v}-\nabla_{\mathbf{x}}\mathcal{V}=0
  \quad\Longrightarrow\quad
  \mathbf{v}^*(\mathbf{x},t)=\tfrac12H^{-1}\nabla_{\mathbf{x}}\mathcal{V}.
\end{equation}
This is the optimal feedback law: the selling rate as a function of the current
inventory and time.

\subsection*{II.4 Substitute back: the time-dependent equation}

\begin{equation}
  \boxed{\;\partial_t\mathcal{V}-\tfrac14(\nabla_{\mathbf{x}}\mathcal{V})^\top H^{-1}
  (\nabla_{\mathbf{x}}\mathcal{V})+\lambda\,\mathbf{x}^\top C\,\mathbf{x}=0\;}
\end{equation}
It runs backward from $t=T$. The terminal condition is $\mathcal{V}(\mathbf{x},T)=0$
for $\mathbf{x}=\mathbf{0}$ and $+\infty$ otherwise; in practice use a penalty
$\tfrac12\mathbf{x}^\top\Phi\,\mathbf{x}$ with $\Phi$ large.

\subsection*{II.5 Check against the known answer (required before going further)}

Try $\mathcal{V}=\tfrac12\mathbf{x}^\top P(t)\,\mathbf{x}$. Then
$\nabla\mathcal{V}=P\mathbf{x}$ and the PDE reduces to the matrix Riccati equation
\begin{equation}
  \dot P=\tfrac12\,P\,H^{-1}P-2\lambda C,\qquad P(T)=\Phi .
\end{equation}
For one asset ($M=1$, $C=\sigma^2$, $H=\tilde\eta$), the resulting optimal
trajectory must equal Almgren and Chriss Eq.~(17):
$x(t)=X\,\sinh(\kappa(T-t))/\sinh(\kappa T)$ with $\kappa=\sqrt{\lambda\sigma^2/\tilde\eta}$.
If it does not, the derivation has an error.

\subsection*{II.6 Status}

\begin{center}
\begin{tabular}{@{}L{5.2cm}L{8.6cm}@{}}
\toprule
\textbf{Item} & \textbf{Status} \\
\midrule
Integral from discrete objective (Part~I) & Derived by hand; algebra reviewed \\
Bellman equation and PDE (II.1--II.4) & Derived by hand; \textbf{not verified numerically} \\
Riccati check against Eq.~(17) (II.5) & \textbf{To do}; gate for everything after \\
Next stage: linearise the PDE (Cole--Hopf with viscosity) & Not started \\
\bottomrule
\end{tabular}
\end{center}

\subsection*{II.7 Assumptions to state in the paper}

\begin{itemize}[leftmargin=1.5em]
  \item Sales only (no buying back), so the spread term is a constant.
  \item Diagonal impact matrices; an antisymmetric permanent-impact part adds a term
        to the optimality condition (Almgren and Chriss, Eq.~31) that is not
        modelled here.
  \item Continuum limit $\tau\to0$, so $\tilde\eta\to\eta$.
  \item The terminal condition $\mathbf{x}(T)=\mathbf{0}$ is a hard constraint and
        is handled by a large penalty $\Phi$, which affects conditioning later.
\end{itemize}

\end{document}